% !TEX program = LuaLaTeX+se
% !TEX encoding = UTF-8 Unicode

\documentclass[12pt, a4paper]{article}
\usepackage[latin, polish]{babel} % ustawienia języka
\usepackage[T1]{fontenc} % przełączenie na fonty EC (256 znaków)
\usepackage[utf8]{luainputenc} % kodowanie
\usepackage{fullpage} % redukcja marginesów
\usepackage{setspace} % odstępy między liniami
\usepackage{geometry} 
\geometry{
	left=12mm,
	right=12mm,
	top=12mm,
	bottom=12mm,
	marginparwidth=0mm,
	marginparsep=0mm
	}

\usepackage{fontspec}
\usepackage{libertine} % ust. czcionki
\pagestyle{empty}	% usuwanie nagłówków/stopek z dokumentu
\usepackage[indent=0pt]{parskip} % ustawianie paragrafów
\setstretch{1.5} 

%  Konfiguracja kolumn
\usepackage{paracol}
\setcolumnwidth{8.3cm,4.3cm}
\setlength{\columnsep}{0.4cm}

\usepackage{gregoriotex}
\gresetlinecolor{gregoriocolor}
\grechangestyle{initial}{\setmainfont{VectisMonumentalW01-Regular}\fontsize{23}{23}\selectfont}
% \grechangedim{spacelinestext}{0.7 cm}{scalable}
% \grechangedim{spaceabovelines}{0.43 cm}{scalable}
% \gresetheadercapture{commentary}{grecommentary}{string}
% \grecommentary{}
\gresetheadercapture{annotation}{greannotation}{string}
% Format annotation above initial
\grechangestyle{annotation}{\small\bfseries}%
%\grechangedim{intersyllablespacenotes}{0.18cm}{scalable}%
%\grechangestyle{commentary}{\it}
%\gresetmodenumbersystem{roman-majuscule}
\usepackage{titlesec} % ustawianie nagłówków
\titlespacing*{\section}{0pt}{0pt}{3mm}	
\titleformat{\section}[hang]{\filcenter\large\scshape\bfseries}{}{0pt}{} % format nagłówków

\newfontface\vectis{VectisMonumentalW01-Regular}
\def\greinitialformat#1{%
{\fontsize{24}{24}\vectis #1}}

\titlespacing*{\section}
{0pt}{0pt}{0pt}

\begin{document}
\large
\begin{spacing}{0.7}
\begin{center}
	\section*{Jesu dulcis memoria}
	\small\textit{Liber Usualis, Solesmes, 1961, p. 452}
	\vspace{15mm}
\end{center}	
\end{spacing}
\greannotation{Hymn.} 
\greannotation{1.} 
\gabcsnippet{QUem(df) ter(e[ll:1]d)ra,(e) pon(g)tus,(g) sí(ef)de(e[ll:1]d)ra(d'_) (,) Co(e)lunt,(g') ad(h)ó(h)rant,(g) praé(hj)di(ivHG)cant,(h.) (;) Tri(g)nam(fe) re(d)gén(e[ll:1]d)tem(c) má(d')chi(d)nam,(e'_[oh:h]) (,) Clau(dh'/ig)strum(h') Ma(d)rí(e)ae(g) bá(fe)ju(d)lat.(e.) (::)}

\vspace{5mm}
\setstretch{1}% odstęp między liniami
1. Jezusa słodkie wspomnienie, Daje duszy pocieszenie: Wszystkie słodycze przechodzi, Kiedy Jezus wśród nas chodzi, wśród nas chodzi.\\
2. Ani wdzięcznych śpiewów pienie, Ani miłe uszom brzmienie, Świętych myśli nie pomnoży Jak ty Jezu Synu Boży!\\
3. Jezu grzeszników nadziejo! Gdy w pokucie łzy swe leją: Kto cię szuka, kto cię prosi, Skarb twej łaski w sobie nosi.\\
4. Żaden język nie opowi: Żaden wyraz nie wysłowi: Samo uczy doświadczanie, Co jest Jezu! twe kochanie.\footnote{M. M. Mioduszewski, Śpiewnik kościelny, Kraków 1838, s. 160–161.}\\
5. Jezu, bądź naszą radością, Który jesteś przyszłą nagrodą, Niech będzie nasza w Tobie chwała na wszystkie wieki wieków.\footnote{tłum. DR}
\end{document}